\documentclass[UTF8,zihao=-4,fontset=windows]{ctexart}
\usepackage[a4paper,margin=24mm]{geometry}
\usepackage{amsmath,amssymb,mathtools}
\usepackage[hidelinks]{hyperref}
\newcommand{\tr}{\operatorname{tr}}
\newcommand{\diag}{\operatorname{diag}}
\newcommand{\norm}[1]{\lVert#1\rVert}
\title{交替罚参数下的二次 ADMM\\一个二维族的精确失稳条件}
\author{}
\date{2026 年 11 月 \quad 私人研究工作稿}
\begin{document}
\maketitle

\begin{abstract}
对二次等式约束问题，固定罚参数的 ADMM 可以收敛，逐次改变罚参数却会改变迭代矩阵的乘积。本文考察两个具有相同谱的二维正定 Hessian，其中一个相对另一个旋转 $45^\circ$，罚参数在共同谱的两个端点之间交替。固定使用这两个端点中的任一个，迭代都收敛。交替迭代则存在唯一的条件数阈值，由一个三次方程精确定义，数值约为 $126.182$。此外，给出一个由参数变化量保证稳定的充分条件。该例用于区分逐参数稳定性、指定切换序列的稳定性与自适应反馈的稳定性。
\end{abstract}

\noindent\textbf{来源和范围。} 杨书泰的私人云端《动力系统与遍历理论视角下的优化算法研究路线图》第 10 节已记录条件数 $199$ 的有理反例及参数变化预算。本文补出消元步骤，并把固定反例扩展为显式二维族，得到交替端点参数时的阈值。本文使用 AI 辅助整理与推导，待作者核阅，精确阈值的新颖性还需比较文献。

\section{一次 ADMM 迭代怎样变成矩阵映射}
设 $Q,R$ 为对称正定矩阵，考虑
\[
 \min_{x=z}\ \tfrac12 x^\mathsf{T}Qx+\tfrac12 z^\mathsf{T}Rz .
\]
以未按罚参数缩放的 $y$ 表示拉格朗日乘子。一次使用参数 $\rho>0$ 的更新为
\begin{align*}
 x^+&=(Q+\rho I)^{-1}(\rho z-y),\\
 z^+&=(R+\rho I)^{-1}(\rho x^++y),\\
 y^+&=y+\rho(x^+-z^+).
\end{align*}
第二个等式给出 $Rz^+=y+\rho(x^+-z^+)=y^+$。因此无论下一步将 $\rho$ 调为何值，从第一个完整更新之后一直有 $y=Rz$。也可直接用 $y_0=Rz_0$ 初始化，把等式从第一步就保持下去。代入并使用 $(Q+\rho I)^{-1}(Q+\rho I)=I$，可得
\[
 \rho x^++Rz=(Q+\rho I)^{-1}
 \bigl[\rho^2 I-\rho R+(Q+\rho I)R\bigr]z
 =(Q+\rho I)^{-1}(\rho^2 I+QR)z.
\]
于是约化更新为
\begin{equation}
 z^+=M_\rho z,\qquad M_\rho=(R+\rho I)^{-1}(Q+\rho I)^{-1}(\rho^2 I+QR).
 \label{eq:map}
\end{equation}
矩阵乘法的顺序来自两块更新的先后。在一般情况下 $QR\ne RQ$，因此改变次序可能改变轨迹。

\section{同谱旋转例的交替阈值}
取 $\kappa\ge1$，并设
\[
 Q=\diag(1,\kappa),\qquad
 R=\tfrac12\begin{pmatrix}1+\kappa&\kappa-1\\ \kappa-1&1+\kappa\end{pmatrix}.
\]
$R$ 的特征向量分别沿 $(1,-1)$ 和 $(1,1)$，对应特征值 $1$ 和 $\kappa$。两块 Hessian 的条件数同为 $\kappa$。使用式 \eqref{eq:map} 的矩阵对 $\rho=1,\kappa$ 代入，可得两个矩阵各自的特征值为
\[
 \frac12,\qquad s=\frac{3\kappa^2+2\kappa+3}{4(\kappa+1)^2}
              =\frac34-\frac{\kappa}{(\kappa+1)^2}.
\]
验证这项计算，只需由式 \eqref{eq:map} 分别取迹与行列式。对两个参数所得结果相同，迹为 $1/2+s$，行列式为 $s/2$。故这两种固定参数的轨道都线性收敛。

按 $y_0=Rz_0$ 初始化，每轮先用 $1$、再用 $\kappa$。令 $P=M_\kappa M_1$，则 $z_{2j}=P^jz_0$，$z_{2j+1}=M_1P^jz_0$。用 $T=\tr(P)$、$D=\det(P)$ 记其迹与行列式。由 $2$ 阶矩阵逆的显式公式以及行列式的乘法性，
\begin{align}
 T &=\frac{-\kappa^4+52\kappa^3+26\kappa^2+52\kappa-1}
              {64\kappa(\kappa+1)^2}, \notag\\
 D &=\frac{(3\kappa^2+2\kappa+3)^2}{64(\kappa+1)^4}. \label{eq:td}
\end{align}
附录给出两次更新的矩阵元素，可以逐项核对迹。令 $p(\lambda)=\lambda^2-T\lambda+D$。显然 $0<D<1$。式 \eqref{eq:td} 还给出
\[
 p(1)=
 \frac{(\kappa^2+6\kappa+1)(\kappa^4+17\kappa^3+36\kappa^2+17\kappa+1)}
 {64\kappa(\kappa+1)^4}>0.
\]
对实系数二次式，在 $0<D<1$ 下，两个根都在单位圆内当且仅当 $p(1)>0$ 且 $p(-1)>0$。如果根是复共轭，其模为 $\sqrt D$。如果根是实数，它们同号。此时若一个根至少为 $1$，另一个正根小于 $1$，于是 $p(1)\le0$。若一个根至多为 $-1$，另一个负根大于 $-1$，于是 $p(-1)\le0$。反过来，两个实根都在 $(-1,1)$ 时，两个测试值为正。

决定阈值的是 $p(-1)=1+T+D$。通分后的分子是六次回文多项式，除以 $\kappa^3$ 后可用 $\kappa^j+\kappa^{-j}$（$j=1,2,3$）表示。利用 $(\kappa+\kappa^{-1})^2-2=\kappa^2+\kappa^{-2}$、$(\kappa+\kappa^{-1})^3-3(\kappa+\kappa^{-1})=\kappa^3+\kappa^{-3}$，得到
\[
 p(-1)=\frac{-\kappa^3}{64\kappa(\kappa+1)^4}
 \bigl(x^3-123x^2-400x-316\bigr),\qquad x=\kappa+\kappa^{-1}.
\]
记括号内多项式为 $F(x)$。当 $x>0$ 时，$F(x)/x^2=x-123-400/x-316/x^2$ 严格递增，从 $-\infty$ 增至 $+\infty$。因此它只过零一次，且 $F(2)=-1600$。令 $x_*$ 为这个零点，并定义
\[
 \kappa_*=\frac{x_*+\sqrt{x_*^2-4}}2.
\]
因为 $\kappa+\kappa^{-1}$ 在 $\kappa\ge1$ 时严格递增，$\kappa<\kappa_*$ 时交替迭代线性收敛。$\kappa=\kappa_*$ 时两个根为 $-1$ 与 $-D$。沿 $-1$ 的特征方向，偶数步子序列不趋于零。$\kappa>\kappa_*$ 时有 $p(-1)<0$，而 $p(\lambda)$ 在 $\lambda\to-\infty$ 时趋于 $+\infty$。故存在一个实根 $\lambda<-1$，取其特征向量为 $z_0$ 并置 $y_0=Rz_0$，即可得到真实的 ADMM 发散轨道。$x_0$ 无需配合选择，因为第一步直接更新 $x^+$。若实现以 $y/\rho$ 保存乘子，改变罚参数时须同步改变这一缩放量，使未缩放乘子 $y$ 保持同一个值。

以整数端点试算即知 $126<\kappa_*<127$。原云端例的 $\kappa=199$ 给出
\[
 T=-\frac{14467601}{6368000},\qquad
 D=\frac{888099601}{6400000000},\qquad
 p(-1)=-\frac{1443188379401}{1273600000000}.
\]
还可以直接使用有理不等式 $T<-2$、$0<D<1$，它们立即给出 $p(-1)<0$。因此有理例的结论无需计算特征值的小数近似。

\section{参数变化量怎样进入收敛估计}
设 $z_{k+1}=M_{\rho_k}z_k$，所有参数取自同一个集合 $\mathcal I\subset(0,\infty)$。假设存在一族正定矩阵 $P_\rho$，使对应范数 $V_\rho(z)=(z^\mathsf{T}P_\rho z)^{1/2}$ 满足
\[
 M_\rho^\mathsf{T}P_\rho M_\rho\preceq r^2P_\rho,\quad
 P_\sigma\preceq e^{L|\log(\sigma/\rho)|}P_\rho,\quad
 0<\underline pI\preceq P_\rho\preceq\overline pI
\]
其中 $0<r<1$、$L\ge0$，三个矩阵不等式对 $\rho,\sigma\in\mathcal I$ 一致成立。第一式衡量一次更新，第二式衡量更新后把旧范数换成新范数的代价。两者合在一起给出
\[
 V_{\rho_{k+1}}(z_{k+1})
 \le r\,e^{(L/2)|\log(\rho_{k+1}/\rho_k)|}V_{\rho_k}(z_k).
\]
若任意长度 $n$ 的窗口内，$\sum_{k=m}^{m+n-1}|\log(\rho_{k+1}/\rho_k)|\le vn+B$，逐项相乘并利用两个范数的等价常数，得到
\[
 \norm{z_{m+n}}\le
 \sqrt{\overline p/\underline p}\,e^{LB/2}
 (r e^{Lv/2})^n\norm{z_m}.
\]
只要 $r e^{Lv/2}<1$，这个界就允许持续变化的罚参数。该条件依赖整个允许的参数族及变化速度。它也指出了下一步要证明什么：对实际的反馈规则导出变化预算，或直接控制更新矩阵的乘积。

在本文的正定二次模型上，可以具体给出一组常数。设 $\mathcal I=[\rho_-,\rho_+]$ 是正紧区间，记 $A_\rho=R+\rho I$，以及
\[
 C_Q=(\rho I-Q)(\rho I+Q)^{-1},\qquad
 C_R=(\rho I-R)(\rho I+R)^{-1}.
\]
由式 \eqref{eq:map}，
\[
 A_\rho M_\rho A_\rho^{-1}=(I+C_Q C_R)/2 .
\]
展开 $(Q+\rho I)(R+\rho I)+(\rho I-Q)(\rho I-R)$，得到 $2(QR+\rho^2 I)$，即得上式。令 $q_Q$ 是 $\mathcal I$ 上 $\norm{C_Q}$ 的最大值，$q_R$ 类同，两者严格小于 $1$。故 $P_\rho=A_\rho^2$ 给出 $r=(1+q_Qq_R)/2<1$。对 $R$ 的任一特征值 $\lambda>0$，有 $(\lambda+\sigma)/(\lambda+\rho)\le\max\{1,\sigma/\rho\}$，所以可取 $L=2$。两个范数等价常数取 $\underline p=(\lambda_{\min}(R)+\rho_-)^2$、$\overline p=(\lambda_{\max}(R)+\rho_+)^2$。于是 $r e^v<1$ 是这组可算常数给出的充分条件。固定参数的收缩因子和缩放已有系统分析，本节使用的是相应的反射算子估计。\cite{ghadimi,gb}
Xu 等在谱罚参数 ADMM 的分析中讨论了控制罚参数变化的充分条件，包括经过有限步后关闭自适应这一实现方式。Lorenz 与 Tran-Dinh 用正上下界和可和增量分析非平稳 DR/ADMM。本文的窗口预算允许非零的长期平均变化量，代价是要单独验证收缩和范数更换的常数。\cite{xu,lorenz}
本例用预先指定的二周期罚参数序列验证失稳。它不回答某个实际反馈能否生成这条周期，也不刻画给定 $\kappa$ 时所有切换序列的最坏因子。

\appendix
\section{迹的逐项核对}
写 $k=\kappa$。把式 \eqref{eq:map} 的两个逆显式代入，得到
\begin{align*}
 M_1 &=\frac{1}{16(k+1)^2}
 \begin{pmatrix}
 -k^3+11k^2+13k+9 & -(k-1)^3\\
 (k-1)^2(k+3) & k^3+9k^2+11k+11
 \end{pmatrix},\\[-1mm]
 M_k &=\frac{1}{16k(k+1)^2}
 \begin{pmatrix}
 11k^3+11k^2+9k+1 & -(k-1)^2(3k+1)\\
 -(k-1)^3 & 9k^3+13k^2+11k-1
 \end{pmatrix}.
\end{align*}
计算 $\tr(M_kM_1)$ 只需交叉对应的四个乘积。复算材料也直接从式 \eqref{eq:map} 计算两个矩阵，然后核对式 \eqref{eq:td}、阈值多项式及 $\kappa=199$ 的特例。有限有理算术核验与 ADMM 收敛论证分开记录。

\begingroup\small
\begin{thebibliography}{9}
\bibitem{ghadimi} E. Ghadimi, A. Teixeira, I. Shames and M. Johansson. Optimal parameter selection for the alternating direction method of multipliers (ADMM): Quadratic problems. \textit{IEEE Transactions on Automatic Control}, 2015. \href{https://arxiv.org/abs/1306.2454}{arXiv:1306.2454}.
\bibitem{gb} P. Giselsson and S. Boyd. Linear convergence and metric selection for Douglas--Rachford splitting and ADMM. \textit{IEEE Transactions on Automatic Control}, 2017. \href{https://web.stanford.edu/~boyd/papers/pdf/metric_select_DR_ADMM.pdf}{论文}.
\bibitem{xu} Z. Xu, M. Figueiredo and T. Goldstein. Adaptive ADMM with spectral penalty parameter selection. \textit{AISTATS}, 2017. \href{https://proceedings.mlr.press/v54/xu17a.html}{论文}.
\bibitem{lorenz} D. A. Lorenz and Q. Tran-Dinh. Non-stationary Douglas--Rachford and alternating direction method of multipliers: Adaptive stepsizes and convergence. \textit{Computational Optimization and Applications}, 2019. \href{https://arxiv.org/abs/1801.03765}{arXiv:1801.03765}.
\end{thebibliography}
\endgroup
\end{document}
